\section{Results}
\label{sec:results}

All values are medians over three repetitions; the shaded bands in the figures are
the inter-quartile range. Figures share one encoding: colour gives the concurrency
family (blue for threaded, orange for \code{poll}, green for \code{epoll}), and line
style and marker give the socket topology (solid for simplex, dashed for duplex,
dotted for split). Values are comparable only within one scenario and one swept value.

%-----------------------------------------------------------------------------
\subsection{S1: fan-in}
\label{sec:res-fanin}

\begin{figure}[h]
\centering
\includegraphics[width=\linewidth]{s1_fanin}
\caption{S1. Each peer offers 100 msg/s. The node has two CPUs (200\%).}
\label{fig:s1}
\end{figure}

\finding{Thread-per-connection collapses at 1,024 peers.} Table~\ref{tab:s1-1024}
shows the result. \mode{threads}, \mode{pool-sender} and \mode{peer-threads} deliver
10.4--20.8k of the 102.4k msg/s offered. Their median RTT is 16--19\,s, they run about
1,030 threads, and they use both cores. Every multiplexed mode delivers the whole load
with no loss. Each runnable reader thread competes for two CPUs, and the scheduler
cannot keep up.

\begin{table}[h]
\centering\small
\begin{tabular}{@{}lrrrrrrr@{}}
\tablebody{tables/s1_1024.tex}
\end{tabular}
\caption{S1 at 1,024 active peers (102.4k msg/s offered).}
\label{tab:s1-1024}
\end{table}

\finding{Duplex halves the descriptors.} At 1,024 peers the duplex and split modes
use about 1,031 fds, against 2,055 for \mode{poll}/\mode{epoll}. \mode{peer-threads}
also uses two per peer, because each link's pool has an eventfd.

\finding{Below the ceiling, threads are the cheapest mode.} At 256 peers
(Table~\ref{tab:s1-256}), \mode{threads}, \mode{pool-sender} and \mode{peer-threads}
use 49--54\% CPU, with a p99 of 6.8--7.2\,ms. \mode{poll} and \mode{epoll} use
114--132\% for a p99 of about 20\,ms. The duplex and split modes fall in between.
A blocked reader costs nothing, and the work of the active ones spreads over both
cores. A single I/O thread makes at least one system call per ready socket, and all
of them run on one core.

\begin{table}[h]
\centering\small
\begin{tabular}{@{}lrrrr@{}}
\tablebody{tables/s1_256.tex}
\end{tabular}
\caption{S1 at 256 active peers (25.6k msg/s offered; every mode delivers it).}
\label{tab:s1-256}
\end{table}

\finding{An unexplained ordering at 1,024 peers.} The poll-based duplex and split
modes have the lowest latency: p50 60--76\,ms and p99 126--168\,ms. Their epoll twins
are at 147--159\,ms p50 and 654--721\,ms p99. The ordering holds in all three
repetitions. The data does not explain it. One plausible mechanism is the order in
which each pool serves ready sockets: \code{poll}'s round-robin cursor against
\code{epoll}'s ready list. We report this as an observation, not a finding.

%-----------------------------------------------------------------------------
\subsection{S2: idle connections}

\begin{figure}[h]
\centering
\includegraphics[width=\linewidth]{s2_idle}
\caption{S2. Eight active peers at 500 msg/s each, plus idle peers.}
\label{fig:s2}
\end{figure}

\finding{\code{poll()} pays for idle sockets; \code{epoll} does not.}
Figure~\ref{fig:s2} shows the cost. From 0 to 1,000 idle peers:
\begin{itemize}
  \item \mode{poll}'s p99 rises from 1.05 to 3.28\,ms and its CPU from 19\% to 100\%,
        because every wake-up rebuilds and scans the whole \code{pollfd} array;
  \item the poll duplex and split modes have half as many sockets to scan and pay
        about half (about 65\% CPU);
  \item every \code{epoll} variant stays flat, at about 1.0\,ms p99 and 10--17\% CPU.
\end{itemize}
In the thread modes each idle peer costs a thread (1,016 threads), but latency stays
flat: a blocked thread uses no CPU.

%-----------------------------------------------------------------------------
\subsection{S3: slow consumers}
\label{sec:res-slow}

\begin{figure}[h]
\centering
\includegraphics[width=0.55\linewidth]{s3_slow}
\caption{S3. The p99 of all peers when some stop reading.}
\label{fig:s3}
\end{figure}

\finding{Only a blocking sender lets one slow reader stall everyone.} One slow peer
raises \mode{threads}' p99 for \emph{all} peers from 1.3\,ms to 419\,ms, and four slow
peers raise it to 2.2\,s (Figure~\ref{fig:s3}). The sender thread waits up to
500\,ms per blocking write. Every other mode stays at 1.3--1.6\,ms. The pooled simplex
modes queue about 2,045 writes per slow peer but never reach the 1\,MiB cap within
20\,s, so nothing is dropped.

\textbf{Caveat.} The duplex, split and \mode{peer-threads} modes record no queued
writes at all, even with a 4\,KiB receive buffer on the slow peers. For these shapes
the result shows that no peer stalled, not that a slow reader was isolated under
backpressure.

%-----------------------------------------------------------------------------
\subsection{S4: churn}
\label{sec:res-churn}

\begin{figure}[h]
\centering
\includegraphics[width=\linewidth]{s4_churn}
\caption{S4. Left: p99 against churn. Right: each valid repetition at churn 10 and 50,
         plotting the node's send failures against p99.}
\label{fig:s4}
\end{figure}

\finding{Churn moves the tail, and the tail follows the sender's failures.} At churn
of 10 reconnects/s and above, most modes reach a p99 of 44--143\,ms. Only
\mode{threads} stays at about 4.5\,ms, and \mode{poll-duplex} does too at churn 10.
Table~\ref{tab:s4} and Figure~\ref{fig:s4} show that p99 follows the number of failed
sends. A failed send (\code{ConnectionNotFound} for a handle that has just been
closed) makes the sender sleep 100\,ms before retrying, on the single thread that
writes for every peer. Every other peer waits with it:
\begin{itemize}
  \item cells with zero or one failure keep a p99 of about 5\,ms;
  \item a handful of failures already raises p99 in some repetitions, and from about
        ten failures per run p99 sits at or above the 100\,ms sleep;
  \item the simplex \mode{poll} and \mode{epoll} modes fail most often (125--166 times
        at churn 50), and they are the only modes whose \emph{median} also rises, to
        32\,ms and 13\,ms.
\end{itemize}

\begin{table}[h]
\centering\small
\begin{tabular}{@{}lllrrr@{}}
\tablebody{tables/s4.tex}
\end{tabular}
\caption{S4: the node's send failures for each of the three repetitions, against tail latency.}
\label{tab:s4}
\end{table}

\finding{Threaded reads do not prevent this.} \mode{pool-sender} and
\mode{peer-threads} read on their own threads, with isolated descriptors, and still
reach a p99 of 64--105\,ms. \mode{threads} has no failures at all. The likely
mechanism, which we have not proven, is that its blocking \code{send()} on a
just-closed socket still succeeds into the kernel buffer. The echo is then lost
silently instead of retried, which also gives \mode{threads} the lowest loss. The pool
modes notice the EOF first, so their next send fails and the sender sleeps. The split
modes fail more often than the duplex modes (69--103 against 14--25 at churn 50). The
same 100\,ms signature appears as the handshake p99 at RTT 0 in S4 and S6a.

%-----------------------------------------------------------------------------
\subsection{S5: throughput ceiling}

\begin{figure}[h]
\centering
\includegraphics[width=\linewidth]{s5_throughput}
\caption{S5. Closed loop, 8 peers with 16 messages outstanding each. Black whiskers are the IQR.}
\label{fig:s5}
\end{figure}

\begin{table}[h]
\centering\small
\begin{tabular}{@{}lrrrr@{}}
\tablebody{tables/s5.tex}
\end{tabular}
\caption{S5: throughput in thousands of messages per second.}
\label{tab:s5}
\end{table}

\finding{The single I/O thread is the ceiling.} At 4 and 64\,KiB, every mode with one
I/O thread plateaus at about 122--144\% CPU and 51--56k msg/s (4\,KiB), or
7.6--8.6k msg/s (64\,KiB). The thread modes reach 170--180\% and the highest
throughput:
\begin{itemize}
  \item 105--113k msg/s at 4\,KiB;
  \item 17.7k msg/s for \mode{pool-sender} and 16.2k for \mode{peer-threads} at
        64\,KiB.
\end{itemize}

\finding{A separate write thread pays off once that thread is the bottleneck.}
Compared with their duplex twins, the split modes gain:
\begin{itemize}
  \item at 4\,KiB, with only 8 peers: 46\% (\mode{poll}) and 76\% (\mode{epoll});
  \item at 64\,KiB: 21\% and 40\%;
  \item at 32\,B: 15\% and 9\%.
\end{itemize}

\finding{Simplex buys no throughput.} At 32\,B, \mode{poll} and \mode{epoll} are
slower than their duplex versions: 158k against 188k, and 169k against 189k msg/s. At
4 and 64\,KiB the two shapes are within noise.

%-----------------------------------------------------------------------------
\subsection{S6: WAN churn and bottleneck}
\label{sec:res-wan}

\begin{figure}[h]
\centering
\includegraphics[width=\linewidth]{s6a_wan_churn}
\caption{S6a. Left: handshake p50 against RTT, with 2\,RTT and 3\,RTT guides. Right:
         delivered rate. Crosses mark the simplex cells that collapsed and were
         excluded as degraded.}
\label{fig:s6a}
\end{figure}

\finding{A handshake costs 3\,RTT in simplex and 2\,RTT in duplex.} At 20\,ms RTT the
simplex modes take 62\,ms to become ready and the duplex modes 41\,ms. At 50\,ms the
duplex modes take about 102\,ms (Figure~\ref{fig:s6a}).

\finding{The simplex collapse is a saturation cliff.} In every simplex mode, the
dial-back is a blocking \code{connect()} on the single sender thread. A churn event
costs that thread about two connects, so about $2\cdot\mathrm{RTT}$ of service time.
Its utilisation is $\rho = \mathrm{churn}\times 2\cdot\mathrm{RTT}$:
\begin{itemize}
  \item at 20\,ms, $\rho = 0.64$: stable, although the initial burst of 32 serialised
        connects gives a handshake p99 of 654\,ms;
  \item at 30\,ms, $\rho = 0.96$: \mode{poll} and \mode{epoll} degrade (2.8k msg/s,
        p50 about 100\,ms), while \mode{threads} and \mode{pool-sender} still hold;
  \item at 40 and 50\,ms, $\rho > 1$: all four simplex modes collapse. Only 10--13 of
        32 peers are ready, and the gates exclude all 24 cells.
\end{itemize}
The duplex modes never dial back and stay at about 3.1k msg/s throughout.

S6c repeats RTT 50\,ms with the churn halved to 8/s ($\rho = 0.8$). The simplex modes
recover to 3.0k msg/s, with a handshake p50 of 155\,ms (3\,RTT) and a p99 of 1.7\,s.
The collapse is therefore a cliff at $\rho = 1$, not a gradual slope.

\finding{What causes the cliff, and what defines it.} The sender's 100\,ms retry
sleep, which sets the tail in S4, plays no part here. The collapsed cells record only
1--4 failed sends per run, which is at most 0.4\,s of sleep in a 25\,s run. The cause
is the serial, blocking dial-back. The \emph{sharpness} of the cliff, however, also
comes from how the benchmark defines churn and validity:
\begin{itemize}
  \item Churn is open loop and always closes a \emph{ready} peer
        (Section~\ref{sec:gen}). Once the sender completes handshakes more slowly
        than churn removes ready peers, the set of ready peers drains and the backlog
        of dial-backs only grows. In the collapsed cells, only 278--341 of the 400
        scheduled churn events found a ready peer to close, the handshake p50 grew to
        2.6--3.2\,s, and the generator received no echoes in the window.
  \item A cell is excluded once fewer than 90\% of its peers are ready at the end.
\end{itemize}
A churn schedule that waits for each redial to finish, or a longer window, would
likely show the same overload as a gradual slowdown. The limit itself, one blocking
dial-back per churn event on one shared thread, belongs to the node. Where along that
curve a cell is called ``collapsed'' belongs to the benchmark.

\begin{table}[h]
\centering\small
\begin{tabular}{@{}lrrrrr@{}}
\tablebody{tables/s6b.tex}
\end{tabular}
\caption{S6b: msg/s at three node egress rates (16\,KiB messages, 4\,ms RTT), queued
         sends at 50\,Mbit/s, and handshake p50 at 1\,Gbit/s.}
\label{tab:s6b}
\end{table}

\finding{Behind a bottleneck, the link sets the pace.} In S6b every mode reaches the
same link-limited rate (Table~\ref{tab:s6b}), about 95\% of line rate. Writes queue as
designed: 4--5k queued sends at 50\,Mbit/s, and fewer than 60 at 1\,Gbit/s. Split and
duplex are identical within noise, so a separate write thread gains nothing when the
link is the limit. Simplex handshakes take about 78\,ms against 9\,ms, because the
dial-back's Ack queues behind data on the rate-limited node egress.

%-----------------------------------------------------------------------------
\subsection{S7: capped fan-in}

\begin{figure}[h]
\centering
\includegraphics[width=\linewidth]{s7_capped}
\caption{S7. 1,024 peers offered against a node cap. 1040 is effectively uncapped.}
\label{fig:s7}
\end{figure}

\begin{table}[h]
\centering\small
\begin{tabular}{@{}lrrrrrr@{}}
\tablebody{tables/s7_256.tex}
\end{tabular}
\caption{S7 at \code{-{}-max-peers 256}.}
\label{tab:s7}
\end{table}

\finding{A cap turns collapse into shedding.} With a cap of 256
(Table~\ref{tab:s7}), \mode{threads} serves its 256 peers with no loss, a p99 of
7.1\,ms and 50\% CPU. Uncapped, it delivers 15.6k msg/s and loses 1.7\,M messages.
\mode{peer-threads} reaches a p99 of 6.8\,ms at 58\% CPU. In both, evictions equal
$1024 - \mathrm{cap}$ exactly. Capped \mode{epoll} uses 115\% CPU, the same as S1
\mode{epoll} at 256 peers (114\%). That is simply the cost of a single I/O thread
serving 256 peers (Section~\ref{sec:res-fanin}), not re-dial churn.

\textbf{The cap is not uniform across modes:}
\begin{itemize}
  \item the duplex and split pools hold \code{max\_peers}${}+16$ connections, so they
        serve 80, 144 and 272 peers;
  \item the dual pools hold $2\times$\code{max\_peers} connections, shared with
        dial-backs, so they serve the cap or one fewer;
  \item the thread modes serve exactly the cap.
\end{itemize}
Multiplexed prunes are also not counted as evictions. Modes must therefore be compared
by the number of peers they served, not by the nominal cap.
